\section{Heights, finite point lists, and rational roots}
\label{sec:heights}

\subsection{An exact arithmetic form of Northcott finiteness over \texorpdfstring{$\Q$}{Q}}

For a canonical primitive vector $x$, define the multiplicative height
\[
                 \Ht([x_0:\cdots:x_n])=\max_i |x_i|.
\]
It is a positive integer. For $r=a/b$ in lowest terms with $b>0$,
$\Ht(r)=\max(|a|,b)$, by identifying $r$ with $[a:b]$.
We use no logarithms or real-valued heights.

\begin{theorem}[Bounded-height enumeration]\label{thm:northcott}
There is a primitive recursive algorithm $\mathcal L(n,B,F_1,\ldots,F_s)$
whose output is a duplicate-free list of exactly the points
\[
 \{P\in\PP^n(\Q):\Ht(P)\le B,\ F_1(P)=\cdots=F_s(P)=0\},
\]
for homogeneous integral input polynomials. Its length is at most
$(2B+1)^{n+1}$. Its correctness is provable in $\iPRA$.
\end{theorem}
\begin{proof}
Visit each vector in $[-B,B]^{n+1}$ in a fixed lexicographic order. Retain
it exactly when its gcd is one, its first nonzero entry is positive, and all
polynomial values vanish. There are precisely $(2B+1)^{n+1}$ visits. Every
retained vector is a point of the required height on the locus. Conversely
the canonical vector of any required point belongs to this box and passes
every test. Uniqueness is Lemma~\ref{lem:normalize}.

For $B=0$ the only visited vector is zero, which is rejected. For $B\ge1$,
an explicit numerical implementation uses a visit counter with base
$(2B+1)$ digits of length $n+1$; subtract $B$ from each digit. The output
list and its length are maintained by a bounded recursion. Polynomial
evaluation on the box satisfies
\[
    |F(x)|\le \|F\|_1\max(1,B)^{\deg F},\qquad
    \|F\|_1=\sum_{\alpha}|c_\alpha|,
\]
because each monomial is bounded by the corresponding power. This estimate
also bounds coefficient-weighted monomials and their partial sums. A raw
monomial computation is bounded by $\max(1,B)^{\deg F}$; list codes are
bounded separately by iterating the fixed pairing function for the stated
number of visits. With the input point as an extra parameter,
the invariant ``if its visit has occurred and it passes the tests, it is in
the retained list'' is bounded. This proves completeness using the Boolean
induction of Section~\ref{sec:foundations}.
\end{proof}

The theorem is the rational, fixed-field form of the height finiteness
principle. It supplies a complete list below a \emph{given} height, not an
upper bound for all points of every variety. Those are different assertions.
When $B=0$ the list is empty, as it must be.

\begin{corollary}[Effective use of a global height bound]\label{cor:heightbound}
For a homogeneous projective polynomial system $V$, suppose an integer $B$ and a proof of
$P\in V(\Q)\Rightarrow\Ht(P)\le B$ are supplied. Then
$\mathcal L(n,B,V)$ is the complete rational-point list. In particular,
its emptiness decides whether $V$ has a rational point.
\end{corollary}
\begin{proof}
Soundness is the previous theorem. The supplied implication reduces every
point to its completeness clause. No search for $B$ is part of the argument.
\end{proof}

\subsection{Heights under displayed coordinate maps}

\begin{proposition}[Upper height bounds]\label{prop:upperheight}
Let $F=(F_0,\ldots,F_m)$ be homogeneous forms of degree $d$ and suppose
the evaluated vector is nonzero at $P$. Put
$A=\max_i\|F_i\|_1$. Then
\[
                 \Ht(F(P))\le A\Ht(P)^d.
\]
If an integral square matrix $M$ has nonzero determinant and row norm
$\|M\|_\infty=\max_i\sum_j|M_{ij}|$, then
\[
 \Ht(MP)\le\|M\|_\infty\Ht(P),\qquad
 \Ht(P)\le\|\operatorname{adj}M\|_\infty\Ht(MP).
\]
\end{proposition}
\begin{proof}
The monomial estimate above bounds the unnormalized image coordinates.
Normalization only decreases them. For the inverse bound, apply the first
inequality to the integral adjugate and use
$\operatorname{adj}(M)M=\det(M)I$. Multiplication by the nonzero scalar
$\det(M)$ fixes the projective point. The adjugate entries are finite
determinant sums, so the constants are explicit primitive recursive data.
\end{proof}

For completeness, the adjugate identity has an elementary finite proof.
The diagonal entry is a cofactor expansion of the determinant. An off-diagonal
entry is a determinant with two equal rows, which is zero: pair each
permutation summand with the summand obtained by exchanging the two row
assignments. Their products agree and their signs are opposite. This is a
bounded involution on a list of at most $k!$ permutations for a $k\times k$
matrix, so its proof is available in $\iPRA$ by finite sums. This is also the
determinant fact used in the next section.

\subsection{A complete zero-dimensional example: rational roots}

\begin{theorem}[Rational-root bounds and a complete list]\label{thm:roots}
Let $f(T)=a_0+\cdots+a_dT^d$ be nonzero, with $a_d\ne0$, and let
$k$ be the least index with $a_k\ne0$. Every nonzero rational root $p/q$
in lowest terms, $q>0$, satisfies
\[
                  p\mid a_k,\qquad q\mid a_d.
\]
It therefore has height at most $\max(|a_k|,|a_d|)$. All rational roots
are obtained by the finite divisor search, with $0$ included exactly when
$a_0=0$. This search and its correctness are primitive recursive and
provable in $\iPRA$.
\end{theorem}
\begin{proof}
Remove the factor $T^k$ and write $g(T)=b_0+\cdots+b_eT^e$ with
$b_0=a_k\ne0$ and $b_e=a_d\ne0$. For $e=0$ there is no nonzero root.
For $e>0$, clear denominators in $g(p/q)=0$:
\[
       b_0q^e+b_1pq^{e-1}+\cdots+b_ep^e=0.
\]
All but the first summand are divisible by $p$, so $p\mid b_0q^e$.
Coprimality of $p,q^e$ and Bezout give $p\mid b_0$. Likewise reducing
the equality modulo $q$ yields $q\mid b_ep^e$ and hence $q\mid b_e$.
The divisibility bounds follow because both coefficients are nonzero.
Enumerate $-|b_0|\le p\le|b_0|$, $p\ne0$, and
$1\le q\le|b_e|$, retain coprime pairs with these divisibilities, and
evaluate $f(p/q)$ exactly. Soundness is the test and completeness is the
argument just given. Finally $f(0)=a_0$ decides the zero root separately.
\end{proof}

For example $T^3-5T^2+6T$ has the complete list $0,2,3$. The zero
polynomial is deliberately excluded: every rational would be a root. A
nonzero constant has the empty list. These cases matter for a total
certificate checker.
